%\section{Introduction}

\section{What are the ``general principles/laws" of positive ODEs?}\lbl{s:wha}

Population dynamics,  ecology,  mathematical epidemiology (ME), immuno-virology, rumour spreading,  chemical reaction networks (CRN) theory, are a few of the biological
interaction networks (BIN) subfields, which study all positive dynamical systems, and have similar preoccupations: the existence and multiplicity of equilibria, their local and
global stability, the occurrence of bifurcations,   persistence, permanence, extinction, etc.

\beD[Positive / non-negative ODE]
An ODE
\[
u_t=f(u),
\]
 is called \emph{positive}~\cite{rantzer2015scalable} or
\emph{non-negative}~\cite{haddad2010} if the non-negative orthant
\[
\mathbb{R}^n_{\ge 0} := \{ u \in \mathbb{R}^n : u_i \ge 0,\; i = 1,\dots,n \}
\]
is forward invariant under the flow; equivalently,
\[
u_i=0,\qquad u_j\ge0\ (j\ne i)
\quad\Longrightarrow\quad
f_i(u)\ge0 .
\]
\eeD






Although the  fields enumerated above are aware of each other  -- see  for  example \cite{AndKurtz}, they have grown quite disconnected due to different examples of interest.  For example, ecology has focused a lot on Lotka-Volterra equations
which factor $u_i'= g_i(u) u_i, i=1, ..., u_i \in \R_+$, and mathematical epidemiology (ME) has focused a lot on models with ``strains" characterized by  factorizations $u_i'= M_i(u) u_i, i=1, ...$, where each $u_i$ is a vector, and $M_i(u)$ are matrices; note that in \CRN s (CRNs), the boundary faces $u_i=0$ are called  siphon faces.
The increased complexity of ME models is compensated by the fact that $M_i(u)$ admit regular splittings \cite{Varga,Van}, which allow defining \NGMs, which are quite often of rank one.

It seems fair to say that with the exception of CRN, the other positive ODE fields have evolved historically mainly as collections of examples, with few general principles/laws, the most notable exception being the next generating matrix (NGM) theorem of \cNGM.
Many other theorems have been developed, but they pertain mostly to examples.

The fact that  ecology and epidemiology models  share lots of common points has led to the term eco-epidemiology -- see  for  example \cite{BateHil}, which is justified by the fact that while ecology and epidemiology are separated practically by their preferences of models,
mathematically  they are quite hard to distinguish.

  We make below the point that  for producing general results which apply equally to examples from ecology,  mathematical epidemiology (ME), immuno-virology, rumour spreading, ...,  the adoption of the CRN framework is a possibly useful  step.
  This suggests rigorous definitions (maybe not final) for concepts which have been previously heuristic, like \DFE\ (DFE), epidemiological model and multi-strain model, and transforms isolated results which might have been known essentially to only few researchers, like the Freedman--Ruan-Tang uniform persistence theorem \cite{FreRuan}, or  the Rahman-Zou competitive exclusion
partitions \cite{Rahman} into the foundations of a unified theory for multi-strain
models.

\subsection{Why CRN theory?}
An important characteristic of the  models in mathematical epidemiology (ME),
ecology, and immuno-virology, etc is that  important phenomena like extinction or persistence occur
when trajectories reach the boundary of the positive orthant, and hence that the notion
of forward--invariant coordinate faces is crucial (these appear in ME  as disease--free or
partially disease--free manifolds).

Now, while boundary faces have been touched upon also in
classical ODE theory, the systematic study of phenomena arising for
 positive ODEs was initiated in  CRN works, see notably \cite{hun,Erdi}, who showed that
all positive polynomial models admit (non-unique) mass-action representations.



The CRN framework we use here is encapsulated by the following definition \cite{AdLS}:
\beD[Stoichiometric and chemical representation]\lbl{d:crn}

A stoichiometric representation of $f$ is a pair $(\Gamma,r)$ such that
\[
f(u)\eq\Gamma r(u),
\qquad
r(u)\ge0,
\]
where $\Gamma\in\R^{n\times n_R}$ is constant and
\[
r:\R_{\ge0}^n\to\R_{\ge0}^{n_R}
\]
is locally Lipschitz.

The representation is called \emph{chemical} if, for every reaction
$\rho$ and every species $i$ satisfying
\[
\Gamma_{i\rho}<0,
\]
there exists a nonnegative locally Lipschitz function
$\widetilde r_{\rho i}$ such that
\[
r_\rho(u)
=
u_i\,\widetilde r_{\rho i}(u).
\]

Equivalently,
\[
u_i=0
\quad\Longrightarrow\quad
r_\rho(u)=0
\qquad
\text{whenever }\Gamma_{i\rho}<0.
\]

Furthermore, for every reaction $\rho$, the rate $r_\rho$ depends only
on species consumed by that reaction and on catalyst species which are
not consumed.  In particular, if
\[
\Gamma_{i\rho}<0,
\]
then $r_\rho$ is non-decreasing in $u_i$.

\eeD



The class of chemical ODEs is very general and is in fact a classic favorite in ME.
Note that for computational implementations, we restrict always  to rational models, which we call ``perturbed" mass-action, and define informally as models which become mass-action after a subset of \para s is set to $0$.  This restriction provides the possibility of investigating first the   computationally easier mass-action case.



A crucial development for the CRN-ME interface was the discovery in \cite{AdLS}
that forward--invariant coordinate faces  correspond to the combinatorial object of  \emph{siphons}  in the theory of Petri nets, and that boundary $\omega-$ limit points with interior starting points must lie in siphon  faces.

%%%%%%%%%




